Los resultados de ordenamiento están en las \Cref{fig:sorting_time} y
\Cref{fig:sorting_memory}, y los de multiplicación de matrices en las
\Cref{fig:matrix_time} y \Cref{fig:matrix_memory}. Cada curva es un
algoritmo, y los valores están promediados sobre todas las muestras,
tipos y dominios,Para mantener el informe conciso, se presentan principalmente 
los resúmenes generales. Los gráficos desglosados por tipo de entrada y dominio 
se generaron igualmente y se encuentran disponibles en data/plots/.

\subsection{Ordenamiento}

Lo primero que salta en la \Cref{fig:sorting_time} es que
\texttt{Sort} gana en prácticamente todos los tamaños, y Quick Sort presenta 
un rendimiento cercano. mientras el arreglo no sea muy grande. El
problema aparece en $n=10^7$: para arreglos ya ordenados (ascendentes o
descendentes), Quick Sort simplemente no terminó dentro del límite de
tiempo y esos casos quedaron marcados como omitidos (el detalle está en
\texttt{quicksort\_omitidos.txt}). Con pivote determinista, este es
justo el escenario que dispara el peor caso $O(n^2)$ de la teoría —y
acá se ve de la forma más contundente posible: ni siquiera alcanza a
correr.

Merge Sort, en cambio, se porta bien en todos los tipos de entrada,
siempre $O(n\log n)$, aunque un poco más lento que \texttt{Sort}
por el costo de ir copiando datos durante la fusión. Patience Sort es
el caso más raro de los cuatro: en teoría también es $O(n\log n)$, pero
en la práctica queda notoriamente atrás, sobre todo con arreglos
aleatorios, por todo el trabajo extra de mantener las pilas y fusionar
con un heap de $k$ vías.

En memoria (\Cref{fig:sorting_memory}) las diferencias son bastante
claras: Merge Sort reserva un buffer del mismo tamaño que el arreglo
($n\times\texttt{sizeof(int)}$ bytes), así que es el que más consume.
\texttt{Sort} y Quick Sort ordenan \textit{in-place}, por lo que
su memoria auxiliar es casi nula (más allá de la pila de recursión).
Patience Sort queda en un punto intermedio (nada despreciable, pero
lejos del extremo de Merge Sort) consistente con sus estructuras de
pilas y el heap de fusión.

\subsection{Multiplicación de matrices}

En la \Cref{fig:matrix_time}, Naive crece cúbico como se espera:
pendiente 3 en la escala log-log, sin sorpresas. Strassen
($O(n^{2.807})$) muestra una pendiente algo menor a partir de $n=256$,
lo que es consistente con su mejor complejidad asintótica. Pero para
$n\leq 64$ pasa lo contrario: Naive es más rápido, porque las sumas y
restas extra de Strassen en cada nivel de recursión pesan más que lo
que se ahorra en multiplicaciones,un ejemplo bastante directo de que
una mejor complejidad asintótica no siempre gana cuando la entrada es
chica.

En memoria (\Cref{fig:matrix_memory}) ambos terminan guardando algo
parecido para la matriz resultado, pero Strassen necesita espacio extra
para las submatrices que arma en cada nivel de recursión, y eso se nota
más mientras más grande es $n$. Naive, al ser puramente iterativo, no requiere 
estructuras auxiliares adicionales de tamaño superior al de las matrices, 
por lo que su uso de memoria crece como \(O(n^2)\).

\begin{figure}[htbp]
    \centering
    \includegraphics[width=0.8\textwidth]{../code/sorting/data/plots/resumen_tiempo.png}
    \caption{Tiempo de ejecución promedio de los algoritmos de ordenamiento en función de \( n \).}
    \label{fig:sorting_time}
\end{figure}

\begin{figure}[htbp]
    \centering
    \includegraphics[width=0.8\textwidth]{../code/sorting/data/plots/resumen_memoria.png}
    \caption{Pico promedio del Working Set utilizado durante la ejecución de los algoritmos de ordenamiento.}
    \label{fig:sorting_memory}
\end{figure}

\begin{figure}[htbp]
    \centering
    \includegraphics[width=0.8\textwidth]{../code/matrix_multiplication/data/plots/resumen_tiempo.png}
    \caption{Tiempo de ejecución promedio de los algoritmos de multiplicación de matrices.}
    \label{fig:matrix_time}
\end{figure}

\begin{figure}[htbp]
    \centering
    \includegraphics[width=0.8\textwidth]{../code/matrix_multiplication/data/plots/resumen_memoria.png}
    \caption{Pico promedio del Working Set utilizado durante la ejecución de los algoritmos de multiplicación de matrices.}
    \label{fig:matrix_memory}
\end{figure}

Vale la pena mencionar un par de cosas que no alcanzan a verse en estos
cuatro gráficos resumen: Quick Sort resulta mucho más sensible al tipo
de entrada que Merge Sort o \texttt{Sort} (que se mantienen
estables sin importar si el arreglo viene ordenado, invertido o
aleatorio); y en matrices, la ventaja de Strassen se nota más en
matrices densas con dominio $D_{10}$, donde el costo extra de sumar y
restar submatrices se diluye mejor. El umbral
\texttt{STRASSEN\_THRESHOLD = 64} de la implementación (ver
\texttt{strassen.cpp}) explica directamente por qué para $n\leq 64$ el
comportamiento es idéntico al de Naive: por debajo de ese tamaño,
Strassen \emph{es} Naive.

En general, la complejidad teórica predice bastante bien el
comportamiento asintótico, pero no cuenta toda la historia: la
implementación concreta, cómo se maneja la memoria, la localidad de
los datos y hasta el tamaño exacto de la entrada terminan siendo
igual de importantes en la práctica, sobre todo en el rango de
tamaños intermedios donde ninguna curva ha "despegado" todavía.